% BANKS-SOURCE: pdf=253 print=243 kind=implicit-exercise id=I-CH11-001
\begin{exercise}[Euclidean sections and Lorentzian correlators]{I-CH11-001}
Banks states that a quantum field theory on a curved space-time is especially
accessible when the complexified space-time has a real Euclidean section.  Make
this statement precise for a real scalar field on a fixed background.

Let a complex manifold $M_{\mathbb C}$ carry a holomorphic metric
$g_{\mathbb C}$.  Assume that it contains a Lorentzian real slice $M_L$ and a
second real slice $M_E$ on which the metric is definite.  In a common static
patch, take the complex metric to have the form
\[
 \dd s_{\mathbb C}^{2}
 =-N(\mathbf x)^2\dd z^2+h_{ij}(\mathbf x)\dd x^i\dd x^j,
 \qquad N>0,
\]
where $h_{ij}$ is positive definite.  Work in $D$ spacetime dimensions, with
$i,j=1,\ldots,D-1$.  The Lorentzian slice has $z=t\in\mathbb R$, while the
Euclidean slice has $z=-\ii\tau$ with $\tau\in\mathbb R$.  Use the mostly-plus
signature $\eta_{\mu\nu}=\operatorname{diag}(-,+,\ldots,+)$ on $M_L$, and
define the positive-definite Euclidean metric by
$g_E=g_{\mathbb C}|_{M_E}$.  Take
\[
 S_L=\int\dd t\,\dd^{D-1}\mathbf x\,N\sqrt h\left[
 \frac{(\partial_t\phi)^2}{2N^2}
 -\frac12h^{ij}\partial_i\phi\partial_j\phi
 -\frac12m^2\phi^2-V_{\mathrm{int}}(\phi)
 \right].
\]

\begin{enumerate}[label=(\alph*)]
 \item Show that $z=-\ii\tau$ makes $g_{\mathbb C}|_{M_E}$ real and
 positive definite, so the stated $g_E$ is positive definite.  Continue the
 action explicitly and prove that
 $\ii S_L\longrightarrow-S_E$.  Write $S_E$ in full.

 \item Define the regulated Euclidean functional integral using Banks's
 continued-source convention $+\ii\int\sqrt{g_E}\,J\phi$, and obtain its
 Schwinger functions by functional differentiation.  State which Euclidean
 boundary condition selects a vacuum and what a smooth Euclidean time circle
 of circumference $\beta$ selects.

 \item Let $H$ generate translations along the static time and set the vacuum
 energy to zero.  Starting with Euclidean-ordered insertions, use the spectral
 representation to derive the analytic continuation to Lorentzian
 time-ordered correlation functions.  Give the continuation with an explicit
 $\ii\epsilon$ prescription for arbitrary time ordering.

 \item State sufficient geometric and field-theoretic conditions for this
 continuation.  Include the coordinate-free role of the two real slices, the
 regularity of the Euclidean geometry, and the conditions that make the
 Euclidean correlators valid boundary data for a Lorentzian theory.

 \item Check the prescription for a free massive scalar in flat space by
 continuing its Euclidean two-point function.  Use
 $p\mathbin{\cdot}x=-p^0t+\mathbf p\mathbin{\cdot}\mathbf x$,
 $p^2=-(p^0)^2+\mathbf p^2$, and the positive-frequency phase
 $\exp(\ii p\mathbin{\cdot}x)$ for $p^0=\omega_{\mathbf p}>0$.  Recover the
 standard positive-frequency Feynman two-point function.
\end{enumerate}
\end{exercise}
